\documentclass[11pt,a4paper]{article}
\usepackage[utf8]{inputenc}
\usepackage[T1]{fontenc}
\usepackage[french]{babel}
\usepackage[margin=2.2cm]{geometry}
\usepackage{graphicx}
\usepackage{booktabs}
\usepackage{tabularx}
\usepackage{array}
\usepackage{xcolor}
\usepackage{listings}
\usepackage{fancyhdr}
\usepackage{caption}
\usepackage{subcaption}
\usepackage{tikz}
\usetikzlibrary{positioning,shapes.multipart}
\usepackage[hidelinks]{hyperref}

\pagestyle{fancy}
\fancyhf{}
\fancyhead[L]{Fatourti}
\fancyhead[R]{Programmation web avancée}
\fancyfoot[C]{\thepage}
\renewcommand{\headrulewidth}{0.4pt}

\captionsetup{font=small,labelformat=empty}
\setlength{\parskip}{0.5em}
\setlength{\parindent}{0pt}

\definecolor{nouveau}{HTML}{006233}
\definecolor{gris}{HTML}{6B6F76}
\lstset{basicstyle=\ttfamily\footnotesize, breaklines=true, frame=single, rulecolor=\color{gris!50},
        columns=fullflexible, keepspaces=true, showstringspaces=false, inputencoding=utf8,
        extendedchars=true, literate={é}{{\'e}}1 {è}{{\`e}}1 {à}{{\`a}}1 {ê}{{\^e}}1 {ç}{{\c c}}1 {°}{{\textdegree}}1}

\newcommand{\capture}[2]{\begin{center}\includegraphics[width=0.92\linewidth]{images/#1}\captionof{figure}{#2}\end{center}}
% deux captures côte à côte, à leur place dans le texte (pas de flottant)
\newcommand{\deuxcaptures}[4]{\par\noindent
  \begin{minipage}[t]{0.49\linewidth}\centering\includegraphics[width=\linewidth]{images/#1}\captionof{figure}{#2}\end{minipage}\hfill
  \begin{minipage}[t]{0.49\linewidth}\centering\includegraphics[width=\linewidth]{images/#3}\captionof{figure}{#4}\end{minipage}\par\medskip}

\begin{document}

% ---------------------------------------------------------------- Page de titre
\begin{center}
Génie Informatique 2\\[0.3em]
Année universitaire 2026/2027\\[2.5em]
{\Huge Mini projet}\\[0.8em]
{\Large Application de gestion des factures}\\[0.2em]
{\large d'eau et d'électricité}\\[2em]
{\Large Rapport final}\\[2.5em]
\textbf{Module :} Programmation web avancée\\[0.4em]
\textbf{Encadrant :} Mohammed Al Achhab\\[2.5em]
\textbf{Réalisé par :}\\[0.4em]
Omar El Horre\\ Othmane Farhane\\ Ilyass Garhani
\end{center}

\section*{Présentation}
Fatourti permet à un fournisseur d'eau et d'électricité de gérer ses clients, leurs compteurs et leurs contrats.
Chaque mois, le client déclare les index de ses compteurs avec une photo de chacun ; la facture est calculée
automatiquement, ou bloquée en anomalie jusqu'à sa vérification par le fournisseur. Un agent dépose une fois
par an le relevé de sa tournée, comparé aux déclarations du client.

Depuis la version alpha, l'application couvre tout le cahier des charges et les consignes ajoutées en cours
(table \texttt{quartier}, plusieurs contrats par client). Elle propose aussi le paiement en ligne, des
notifications par e-mail et en AJAX, et une nouvelle interface pensée comme un service public marocain.

% ---------------------------------------------------------------- Conformité
\section{Conformité au cahier des charges}
\begin{center}\small
\begin{tabularx}{\linewidth}{>{\raggedright\arraybackslash}p{6.2cm}X}
\toprule
\textbf{Exigence} & \textbf{Réalisation} \\
\midrule
Le client consulte ses factures & Page \textit{Mes relevés} : chaque facture (PDF), son état et son échéance. \\
Saisie mensuelle avec photos des compteurs & Formulaire \textit{Déclarer ma consommation} (deux index, deux photos JPEG de 3 Mo max). \\
Index et photo du mois précédent affichés & Affichés au-dessus du formulaire, pour comparer avant de saisir. \\
Facture générée automatiquement sans anomalie & Calcul par tranches + TVA 18 \%, PDF créé à l'enregistrement. \\
Anomalie : index inférieur ou seuil dépassé & Seuils réglables par chaque fournisseur (\textit{Mon compte}). \\
Facture : numéro, date, période, client, consommations, HT, TVA, TTC, photos & PDF de facture avec le détail par tranche. \\
Réclamation : fuite externe, fuite interne, facture, autre & Formulaire et suivi côté client ; traitement et réponse côté fournisseur. \\
Le fournisseur gère ses clients & Ajouter, modifier, supprimer ; plusieurs adresses (compteurs) par client. \\
Saisie ouverte du 18 à la fin du mois & Vérifiée côté serveur, rappelée dans l'interface. \\
Correction d'une anomalie avant la facture & Page de correction avec les photos, puis génération de la facture. \\
Traiter les réclamations et notifier le client & Réponse enregistrée, notification dans l'application et par e-mail. \\
TVA 18 \% ; tranches électricité (0,82 / 0,92 / 1,10) et eau (2,37 / 7,39 / 10,98 / 11,03) & Constantes \texttt{TRANCHES\_ELECTRICITE} et \texttt{TRANCHES\_EAU}, appliquées progressivement. \\
Fichier annuel \texttt{Consommation\_Annuelle.txt} & Import par l'agent, contrôle de chaque ligne (agent, année, date, client). \\
Tableau de bord & Impayés, évolution mensuelle (courbe), anomalies, réclamations, dernières factures. \\
Vérification annuelle par le fournisseur & \textit{Bilan annuel} : déclaré / relevé par l'agent / reste à payer, rappels. \\
Enregistrement du paiement & Paiement en espèces par le fournisseur, ou en ligne par le client. \\
\midrule
\textit{Consignes ajoutées en cours :} table quartier ; plusieurs contrats dans des quartiers différents ; un quartier a un type d'activité &
Table \texttt{quartier} (nom, type d'activité) ; chaque compteur est dans un quartier et a son propre contrat. \\
\bottomrule
\end{tabularx}
\end{center}

% ---------------------------------------------------------------- Ancienne interface
\section{L'interface de la version alpha}
La version alpha reposait sur un menu latéral sombre et des pages larges, conçues pour un écran de 1920 pixels.
Les formulaires ne demandaient que l'essentiel (nom, prénom, une adresse, e-mail) et chaque client n'avait
qu'une seule adresse.

\deuxcaptures{ancien-connexion.jpg}{Connexion}{ancien-clients.jpg}{Liste des clients}
\deuxcaptures{ancien-ajout-client.jpg}{Ajout d'un client : une seule adresse}{ancien-releves-client.jpg}{Déclaration et historique du client}

% ---------------------------------------------------------------- Base de données
\section{Évolution de la base de données}
La base passe de 14 tables (50 colonnes) à 19 tables (91 colonnes). Aucune table ni colonne n'a été retirée :
les scripts de \texttt{migrations/} font passer une base existante d'une version à l'autre.

\subsection*{Tables ajoutées}
\begin{center}\small
\begin{tabularx}{\linewidth}{lX}
\toprule
\textbf{Table} & \textbf{Rôle et colonnes} \\
\midrule
\texttt{quartier} & Un quartier et son type d'activité (résidentiel, commercial, industriel) : \texttt{id, nom, typeActivite}. \\
\texttt{adresseCompteur} & Un compteur à une adresse ; un client peut en avoir plusieurs, dans des quartiers différents :
\texttt{id, idConsommateur, adresse, categorie, idQuartier}. \\
\texttt{contrat} & Un contrat par compteur, avec son PDF et son statut (actif, en résiliation, résilié) :
\texttt{id, dateDebut, cheminGeneration, idConsommateur, idAdresseCompteur, statut}. \\
\texttt{notification} & Notifications persistantes (lues, à traiter, résolues) :
\texttt{id, idUtilisateur, type, titre, detail, lien, aTraiter, idObjet, dateCreation, dateLecture, dateResolution}. \\
\texttt{regularisationAnnuelle} & Écart entre le relevé de l'agent et les déclarations, montant estimé, échéance et rappels :
\texttt{id, idConsommationAnuelle, differenceEau, differenceElectricite, montantEstime, dateLimite, dateDernierEnvoi, nombreEnvois, dateReglement}. \\
\bottomrule
\end{tabularx}
\end{center}

\subsection*{Colonnes ajoutées aux tables existantes}
\begin{center}\small
\begin{tabular}{lll}
\toprule
\textbf{Table} & \textbf{Colonnes} & \textbf{Pour} \\
\midrule
\texttt{consommateur} & \texttt{cin, telephone} & Numéro d'identité et téléphone du client \\
\texttt{declaration} & \texttt{idAdresseCompteur} & Le relevé porte sur un compteur précis \\
\texttt{facture} & \texttt{modePaiement, referencePaiement} & Espèces ou en ligne, numéro de transaction \\
\texttt{fournisseur} & \texttt{seuilEau, seuilElectricite} & Seuils d'anomalie définis par le fournisseur \\
\bottomrule
\end{tabular}
\end{center}

\subsection*{Schéma des relations}
Les tables ajoutées sont en vert.
\begin{center}
\begin{tikzpicture}[
  table/.style={draw, rounded corners=2pt, minimum width=3.1cm, minimum height=0.75cm, font=\small\ttfamily, fill=white},
  ajout/.style={table, draw=nouveau, fill=nouveau!8, text=nouveau!80!black, line width=0.9pt},
  lien/.style={->, >=stealth, gray}]
  \node[table] (util) at (0,0) {utilisateur};
  \node[table] (four) at (-4.2,-1.6) {fournisseur};
  \node[table] (cons) at (0,-1.6) {consommateur};
  \node[table] (agent) at (4.2,-1.6) {agent};
  \node[ajout] (notif) at (4.2,0) {notification};
  \node[ajout] (adr) at (-2.2,-3.3) {adresseCompteur};
  \node[ajout] (quart) at (-6.2,-3.3) {quartier};
  \node[ajout] (contrat) at (-4.2,-5) {contrat};
  \node[table] (decl) at (0,-5) {declaration};
  \node[table] (fact) at (0,-6.7) {facture};
  \node[table] (recl) at (4.2,-3.3) {reclamation};
  \node[table] (annu) at (4.2,-5) {ConsommationAnuelle};
  \node[ajout] (regul) at (4.2,-6.7) {regularisationAnnuelle};
  \draw[lien] (four) -- (util); \draw[lien] (cons) -- (util); \draw[lien] (agent) -- (util);
  \draw[lien] (notif) -- (util); \draw[lien] (cons) -- (four);
  \draw[lien] (adr) -- (cons); \draw[lien] (adr) -- (quart);
  \draw[lien] (contrat) -- (adr); \draw[lien] (contrat.north) -- ++(0,0.4) -| (cons.south west);
  \draw[lien] (decl) -- (adr); \draw[lien] (decl) -- (cons);
  \draw[lien] (fact) -- (decl); \draw[lien] (recl) -- (cons);
  \draw[lien] (annu) -- (cons); \draw[lien] (regul) -- (annu);
\end{tikzpicture}
\captionof{figure}{Principales tables et clés étrangères (la flèche va vers la table référencée). Les sous-types de réclamation (fuite interne, externe, facture, autre) ne sont pas représentés.}
\end{center}

% ---------------------------------------------------------------- Architecture et routage
\section{Organisation et routage}
Le code garde l'organisation MVC du cours : \textbf{IHM} (pages), \textbf{Traitement} (contrôle et règles),
\textbf{BD} (accès aux données). \texttt{index.php} est l'\textbf{unique point d'entrée} : le fichier
\texttt{.htaccess} lui envoie toutes les adresses sauf les fichiers statiques de \texttt{public/} (CSS, JS).
Un fichier PHP appelé directement (\texttt{IHM/login.php}, \texttt{BD/connexion.php}\ldots) répond 404, comme
toute adresse inconnue.

Toutes les adresses suivent la forme \texttt{table/action/id}, sans paramètre après un « ? » :
\begin{center}\footnotesize
\begin{tabularx}{\linewidth}{>{\raggedright\arraybackslash}X>{\raggedright\arraybackslash}p{4.4cm}}
\toprule
\textbf{Adresse} & \textbf{Page} \\
\midrule
\texttt{utilisateur/authentification}, \texttt{utilisateur/profil} & Connexion, Mon compte \\
\texttt{consommateur/all}, \texttt{consommateur/ajouter}, \texttt{consommateur/modifier/9107} & Clients \\
\texttt{contrat/resilier/9107} & Résiliation d'un contrat (en deux temps) \\
\texttt{declaration/all}, \texttt{declaration/ajouter/9613}, \texttt{declaration/corriger/9140} & Relevés, déclaration, correction \\
\texttt{facture/all/impayees}, \texttt{facture/detail/9112}, \texttt{facture/telecharger/9112} & Factures \\
\texttt{reclamation/all/attente}, \texttt{reclamation/detail/9108} & Réclamations \\
\texttt{annuelle/all/2026}, \texttt{agent/importer} & Bilan annuel, import de l'agent \\
\texttt{notification/all}, \texttt{notification/lire/97} & Notifications (AJAX, JSON) \\
\texttt{paiement/payer/9112}, \texttt{paiement/retour/9112} & Paiement en ligne \\
\bottomrule
\end{tabularx}
\end{center}
Les filtres sont dans le chemin (\texttt{facture/all/impayees}) ; la recherche des réclamations est envoyée en
POST sur la même adresse.

% ---------------------------------------------------------------- Nouvelle interface
\section{La nouvelle interface}
L'identité visuelle est celle d'un service public : fond « papier », vert du drapeau pour l'action principale
et ce qui est payé ou validé, rouge pour les retards et les anomalies. Les listes partagent un même tableau :
recherche, tri par colonne et pagination.

\subsection*{Espace fournisseur}
\capture{nouveau-tableau-fournisseur.jpg}{Tableau de bord : montant à encaisser, clients, anomalies, consommation mensuelle, réclamations ouvertes.}
\deuxcaptures{nouveau-clients.jpg}{Clients : compteurs, consommation des six derniers mois, impayés}{nouveau-ajout-client.jpg}{Ajout d'un client : plusieurs adresses, chacune dans un quartier}
\deuxcaptures{nouveau-modifier-client.jpg}{Modification : compteurs et contrats du client}{nouveau-resiliation.jpg}{Résiliation : choix du contrat, relevé final}
\deuxcaptures{nouveau-factures.jpg}{Factures : filtre, recherche, tri}{nouveau-facture-detail.jpg}{Détail d'une facture et paiement en espèces}
\deuxcaptures{nouveau-reclamations-fournisseur.jpg}{Réclamations : statut, type, recherche}{nouveau-bilan-annuel.jpg}{Bilan annuel : déclaré, relevé par l'agent, reste à payer}
\deuxcaptures{nouveau-notifications.jpg}{Notifications chargées en AJAX}{nouveau-anomalies.jpg}{Relevés en anomalie à corriger}

\subsection*{Espace client}
\deuxcaptures{nouveau-tableau-client.jpg}{Tableau de bord du client}{nouveau-releves-client.jpg}{Mes relevés : factures, photos, paiement}
\deuxcaptures{nouveau-declarer.jpg}{Déclaration : dernière photo, nouvelle photo et consommation calculée pendant la saisie}{nouveau-reclamations-client.jpg}{Mes réclamations et leurs réponses}
\deuxcaptures{nouveau-mon-compte.jpg}{Mon compte : compteurs, quartiers, contrats}{nouveau-404.jpg}{Toute adresse inconnue : un grand 404, et le lien de retour}

\subsection*{Espace agent}
\capture{nouveau-agent.jpg}{Agent : import du fichier annuel et liste des clients (identifiant, compteurs, déclaré dans l'année, relevé de la tournée).}

% ---------------------------------------------------------------- Cas particuliers
\section{Cas particuliers traités}
\begin{itemize}
\item \textbf{Un relevé par compteur et par mois}, du 18 à la fin du mois. Le relevé de départ (index initial
  saisi par le fournisseur) ne compte pas.
\item \textbf{Mois oublié.} Le relevé suivant est facturé depuis le dernier index validé ; le seuil d'anomalie
  est multiplié par le nombre de mois couverts, pour ne pas signaler une fausse anomalie.
\item \textbf{Client qui ne déclare jamais.} À la fin de l'année, il paie tout ce que l'agent a relevé :
  ses déclarations comptent pour zéro dans le bilan annuel.
\item \textbf{Résiliation en deux temps.} Le client est prévenu et fait un relevé final, même hors période
  et même s'il a déjà déclaré ce mois-ci ; ce relevé crée la facture finale et clôt le contrat.
  Si ce relevé est en anomalie, le contrat attend sa correction. Un compteur résilié ne reçoit plus de relevé.
\item \textbf{Compteur repris par un nouveau client.} Les contrats sont par compteur : l'ancien client règle
  sa facture finale, le nouveau commence avec son propre relevé de départ.
\item \textbf{Suppression d'un client avec un historique.} Le fournisseur est averti (relevés, factures,
  réclamations) et peut résilier plutôt que supprimer ; la suppression se fait dans une transaction.
\item \textbf{Facture en retard.} Après 30 jours, elle apparaît « en retard de N jours ».
\item \textbf{Paiement en ligne.} Au retour, le serveur redemande le paiement au prestataire et vérifie le
  montant, la devise et le numéro de facture ; un retour forgé ou un montant modifié est refusé ;
  une facture déjà payée ne peut pas l'être deux fois.
\item \textbf{Saisies piégées.} Un e-mail comme \texttt{' OR '1'='1} ou un identifiant modifié dans l'adresse
  reste une simple valeur (injection SQL) ; une adresse avec une apostrophe (« Rue de l'Oued ») s'enregistre
  normalement ; les textes affichés sont échappés (XSS) ; les formulaires sensibles ont un jeton CSRF.
\item \textbf{Date de démonstration.} En mode démo, une barre change la date de toute l'application
  (PHP et MySQL) pour montrer la période du 18, le passage au mois suivant ou la fin d'année.
\end{itemize}

% ---------------------------------------------------------------- Conformité au cours
\section{Lien avec le cours}
\begin{itemize}
\item \textbf{PHP et formulaires} : méthodes GET et POST, superglobales \texttt{\$\_GET}, \texttt{\$\_POST},
  \texttt{\$\_FILES}, \texttt{\$\_SESSION}, contrôle avec \texttt{isset} et \texttt{empty}, inclusion de fichiers.
\item \textbf{Bases de données} : PHP/MySQL avec \texttt{mysqli} ; requêtes préparées ou valeurs protégées.
\item \textbf{MVC} : séparation IHM / Traitement / BD et un routeur unique \texttt{index.php}.
\item \textbf{AJAX} : la cloche des notifications utilise l'objet \texttt{XMLHttpRequest} comme dans le cours ;
  le serveur répond en JSON (\texttt{json\_encode}, en-tête \texttt{application/json}) et la page est mise à jour
  par le DOM (\texttt{createElement}, \texttt{appendChild}, \texttt{replaceChildren}) sans rechargement.
\end{itemize}
\begin{lstlisting}
function requeteAjax(methode, adresse, parametres, quandFini, siErreur) {
    const objetXHR = new XMLHttpRequest();
    objetXHR.open(methode, SITE + adresse, true);              // asynchrone
    objetXHR.onreadystatechange = function () {
        if (objetXHR.readyState === 4) {                       // données arrivées
            if (objetXHR.status === 200) quandFini(JSON.parse(objetXHR.responseText));
            else if (siErreur) siErreur(objetXHR.status);
        }
    };
    if (methode === "POST")
        objetXHR.setRequestHeader("Content-Type", "application/x-www-form-urlencoded");
    objetXHR.send(parametres);
}
\end{lstlisting}

% ---------------------------------------------------------------- Conclusion
\section*{Conclusion}
La version finale de Fatourti couvre l'ensemble du cahier des charges et des consignes ajoutées en cours,
avec une base de données enrichie (cinq tables, sept colonnes), un routage unique en \texttt{table/action/id}
et une interface entièrement revue. L'application tourne avec Docker ; la base complète peut être importée
avec le fichier \texttt{databaseTEST.sql}.

\end{document}
